% Module 11: Qiskit Programming Lab
% Estimated slides: 40 | Duration: ~4 hours (hands-on)
% Key topics: QuantumCircuit API, gate methods, measurement, visualization,
%             parameterized circuits, Sampler/Estimator primitives, Qiskit ML,
%             full pipeline from circuit to hardware result
% Labs: progressive Qiskit exercises; capstone mini-project

%%%%%%%%%%%%%%%%%%%%%%%%%%%%%%%%%%%%%%%%%%%%%%%%%%%%%%%%%%
\begin{frame}[fragile]\frametitle{}
\begin{center}
{\Large Module 11: Qiskit Programming Lab}\\[0.5em]
{\large Hands-On Quantum Computing with Python}
\end{center}
\end{frame}

%%%%%%%%%%%%%%%%%%%%%%%%%%%%%%%%%%%%%%%%%%%%%%%%%%%%%%%%%%%
\begin{frame}[fragile]\frametitle{Qiskit Installation and Setup}
\begin{itemize}
\item Install in a fresh conda environment (recommended):
\end{itemize}
\begin{verbatim}
conda create -n quantum python=3.11
conda activate quantum
pip install qiskit qiskit-aer qiskit-ibm-runtime
pip install qiskit-nature qiskit-machine-learning
pip install matplotlib pylatexenc notebook
\end{verbatim}
\begin{itemize}
\item Verify: \texttt{python -c "import qiskit; print(qiskit.version.VERSION)"}
\item Save IBM Quantum credentials (one time):
\end{itemize}
\begin{verbatim}
from qiskit_ibm_runtime import QiskitRuntimeService
QiskitRuntimeService.save_account(channel="ibm_quantum",
                                  token="YOUR_TOKEN")
\end{verbatim}
\end{frame}

%%%%%%%%%%%%%%%%%%%%%%%%%%%%%%%%%%%%%%%%%%%%%%%%%%%%%%%%%%%
\begin{frame}[fragile]\frametitle{QuantumCircuit: The Core Object}
\begin{itemize}
\item \texttt{QuantumCircuit(n)} creates $n$-qubit register (all initialized to $|0\rangle$)
\item \texttt{QuantumCircuit(n, m)} adds $m$ classical bits for measurement results
\item Key methods: \texttt{.h()}, \texttt{.x()}, \texttt{.cx()}, \texttt{.measure()}, \texttt{.draw()}
\item Circuit composition: \texttt{qc1.compose(qc2)}; append: \texttt{qc.append(gate, qubits)}
\item Barrier: \texttt{qc.barrier()} separates sections visually (no physical effect)
\item Reset: \texttt{qc.reset(qubit)} returns qubit to $|0\rangle$ (destructive measurement)
\end{itemize}
\begin{verbatim}
from qiskit import QuantumCircuit
qc = QuantumCircuit(2, 2)
qc.h(0); qc.cx(0, 1)
qc.measure([0,1], [0,1])
qc.draw('mpl')
\end{verbatim}
\end{frame}

%%%%%%%%%%%%%%%%%%%%%%%%%%%%%%%%%%%%%%%%%%%%%%%%%%%%%%%%%%%
\begin{frame}[fragile]\frametitle{Gate Reference: Common Methods}
\begin{itemize}
\item Single-qubit: \texttt{.x(q)}, \texttt{.y(q)}, \texttt{.z(q)}, \texttt{.h(q)}, \texttt{.s(q)}, \texttt{.t(q)}, \texttt{.sdg(q)}, \texttt{.tdg(q)}
\item Rotations: \texttt{.rx(theta,q)}, \texttt{.ry(theta,q)}, \texttt{.rz(theta,q)}, \texttt{.p(phi,q)}
\item Two-qubit: \texttt{.cx(c,t)}, \texttt{.cz(c,t)}, \texttt{.ch(c,t)}, \texttt{.swap(q1,q2)}, \texttt{.iswap(q1,q2)}
\item Three-qubit: \texttt{.ccx(c1,c2,t)}, \texttt{.cswap(c,q1,q2)}
\item Custom unitary: \texttt{.unitary(matrix, qubits)}
\item Inverse: \texttt{gate.inverse()} for any gate
\end{itemize}
% Suggested diagram: Gate symbol reference card
\end{frame}

%%%%%%%%%%%%%%%%%%%%%%%%%%%%%%%%%%%%%%%%%%%%%%%%%%%%%%%%%%%
\begin{frame}[fragile]\frametitle{Simulation with Qiskit Aer}
\begin{itemize}
\item \texttt{AerSimulator}: all-purpose simulator with shot-based sampling
\item \texttt{StatevectorSimulator}: exact state, no measurement sampling
\item Run workflow: \texttt{transpile} then \texttt{backend.run()} then \texttt{result.get\_counts()}
\end{itemize}
\begin{verbatim}
from qiskit_aer import AerSimulator
from qiskit import transpile
sim = AerSimulator()
qc_t = transpile(qc, sim)
job = sim.run(qc_t, shots=1024)
counts = job.result().get_counts()
from qiskit.visualization import plot_histogram
plot_histogram(counts)
\end{verbatim}
\end{frame}

%%%%%%%%%%%%%%%%%%%%%%%%%%%%%%%%%%%%%%%%%%%%%%%%%%%%%%%%%%%
\begin{frame}[fragile]\frametitle{Statevector: Peek Under the Hood}
\begin{itemize}
\item Statevector simulator gives exact quantum state (no measurement needed)
\item Useful for debugging, understanding what a circuit does
\item \texttt{Statevector.from\_instruction(qc)}: simulate without measurement
\item Visualize as bar chart, Bloch sphere, or text representation
\end{itemize}
\begin{verbatim}
from qiskit.quantum_info import Statevector
qc_no_measure = QuantumCircuit(2)
qc_no_measure.h(0); qc_no_measure.cx(0,1)
sv = Statevector(qc_no_measure)
print(sv)  # Shows (|00>+|11>)/sqrt(2)
sv.draw('latex')  # or 'bloch'
\end{verbatim}
\end{frame}

%%%%%%%%%%%%%%%%%%%%%%%%%%%%%%%%%%%%%%%%%%%%%%%%%%%%%%%%%%%
\begin{frame}[fragile]\frametitle{Parameterized Circuits}
\begin{itemize}
\item Use \texttt{Parameter} or \texttt{ParameterVector} for variable gate angles
\item Essential for VQE, QAOA, quantum machine learning
\item Bind parameters with \texttt{qc.assign\_parameters(\{theta: 1.5\})}
\item Can bind multiple circuits at once for efficient optimization
\end{itemize}
\begin{verbatim}
from qiskit.circuit import ParameterVector
theta = ParameterVector('theta', 4)
qc = QuantumCircuit(2)
qc.ry(theta[0], 0); qc.ry(theta[1], 1)
qc.cx(0, 1)
qc.ry(theta[2], 0); qc.ry(theta[3], 1)
bound = qc.assign_parameters([0.5, 1.0, 0.2, 1.5])
\end{verbatim}
% Suggested diagram: Parameterized circuit diagram with theta labels
\end{frame}

%%%%%%%%%%%%%%%%%%%%%%%%%%%%%%%%%%%%%%%%%%%%%%%%%%%%%%%%%%%
\begin{frame}[fragile]\frametitle{Qiskit Runtime Primitives: Sampler and Estimator}
\begin{itemize}
\item New Qiskit 1.x workflow: use Primitives instead of backend.run()
\item \textbf{SamplerV2}: samples measurement outcomes from circuits; returns bit-string counts
\item \textbf{EstimatorV2}: computes expectation values of observables; used in VQE
\item Handles error mitigation automatically when running on hardware
\item Works identically on simulators and real hardware (unified interface)
\end{itemize}
\begin{verbatim}
from qiskit_aer.primitives import SamplerV2
sampler = SamplerV2()
job = sampler.run([qc], shots=1024)
result = job.result()
pub_result = result[0]
counts = pub_result.data.meas.get_counts()
\end{verbatim}
\end{frame}

%%%%%%%%%%%%%%%%%%%%%%%%%%%%%%%%%%%%%%%%%%%%%%%%%%%%%%%%%%%
\begin{frame}[fragile]\frametitle{Circuit Visualization and Analysis}
\begin{itemize}
\item Draw circuits: \texttt{qc.draw('mpl')} (matplotlib), \texttt{'text'} (ASCII), \texttt{'latex'} (LaTeX)
\item Count gates: \texttt{qc.count\_ops()} --- shows gate histogram
\item Circuit depth: \texttt{qc.depth()}
\item Circuit width: \texttt{qc.num\_qubits}
\item DAG representation: \texttt{circuit\_to\_dag(qc)} for analysis
\item OpenQASM export: \texttt{qc.qasm()} for interoperability with other tools
\end{itemize}
\begin{verbatim}
print(qc.count_ops())  # {'cx': 1, 'h': 1, 'measure': 2}
print(qc.depth())       # 3
print(qc.num_qubits)    # 2
\end{verbatim}
\end{frame}

%%%%%%%%%%%%%%%%%%%%%%%%%%%%%%%%%%%%%%%%%%%%%%%%%%%%%%%%%%%
\begin{frame}[fragile]\frametitle{Lab Exercise 1: GHZ State}
\begin{itemize}
\item Build a 3-qubit GHZ state: $(|000\rangle + |111\rangle)/\sqrt{2}$
\item Steps: H on qubit 0; CNOT(0$\to$1); CNOT(0$\to$2)
\item Measure all 3 qubits; expected: only 000 and 111 outcomes
\item Extension: add a bit-flip error on qubit 1 and observe the effect
\end{itemize}
\begin{verbatim}
qc = QuantumCircuit(3, 3)
qc.h(0)
qc.cx(0, 1)
qc.cx(0, 2)
qc.measure([0,1,2], [0,1,2])
\end{verbatim}
\end{frame}

%%%%%%%%%%%%%%%%%%%%%%%%%%%%%%%%%%%%%%%%%%%%%%%%%%%%%%%%%%%
\begin{frame}[fragile]\frametitle{Lab Exercise 2: Random Circuit}
\begin{itemize}
\item Generate a random quantum circuit and analyze its properties
\item Qiskit's \texttt{random\_circuit} function for quick generation
\item Analyze depth, gate counts; check statevector for entanglement
\end{itemize}
\begin{verbatim}
from qiskit.circuit.random import random_circuit
qc = random_circuit(num_qubits=4, depth=5,
                    measure=True, seed=42)
qc.draw('mpl')
print(qc.depth(), qc.count_ops())
\end{verbatim}
\end{frame}

%%%%%%%%%%%%%%%%%%%%%%%%%%%%%%%%%%%%%%%%%%%%%%%%%%%%%%%%%%%
\begin{frame}[fragile]\frametitle{Lab Exercise 3: Bernstein-Vazirani in Qiskit}
\begin{itemize}
\item Implement Bernstein-Vazirani for hidden string $s = 101$
\item Build oracle circuit for $f(x) = s \cdot x$ using CNOT gates
\item Apply H on all qubits, oracle, H again; measure
\item Expected output: $101$ in binary --- reveals hidden string in 1 query
\end{itemize}
\begin{verbatim}
n = 3; s = '101'
qc = QuantumCircuit(n+1, n)
qc.x(n); qc.h(range(n+1))
for i, bit in enumerate(reversed(s)):
    if bit == '1': qc.cx(i, n)
qc.h(range(n)); qc.measure(range(n), range(n))
\end{verbatim}
\end{frame}

%%%%%%%%%%%%%%%%%%%%%%%%%%%%%%%%%%%%%%%%%%%%%%%%%%%%%%%%%%%
\begin{frame}[fragile]\frametitle{Lab Exercise 4: Noise Model Simulation}
\begin{itemize}
\item Create a Bell state circuit
\item Add a depolarizing noise model (1\% error per gate)
\item Compare ideal vs noisy histogram
\item Observe: 01 and 10 outcomes appear due to noise
\end{itemize}
\begin{verbatim}
from qiskit_aer.noise import NoiseModel
from qiskit_aer.noise.errors import depolarizing_error
noise_model = NoiseModel()
error_1q = depolarizing_error(0.01, 1)
error_2q = depolarizing_error(0.05, 2)
noise_model.add_all_qubit_quantum_error(error_1q, ['h'])
noise_model.add_all_qubit_quantum_error(error_2q, ['cx'])
\end{verbatim}
% Suggested diagram: Ideal vs noisy histogram comparison
\end{frame}

%%%%%%%%%%%%%%%%%%%%%%%%%%%%%%%%%%%%%%%%%%%%%%%%%%%%%%%%%%%
\begin{frame}[fragile]\frametitle{Qiskit Machine Learning: Quantum SVM}
\begin{itemize}
\item Quantum feature map: encode classical data as quantum state
\item ZZFeatureMap: entangling feature map using data-driven ZZ interactions
\item Quantum Kernel: compute kernel matrix using quantum circuit fidelities
\item Quantum SVM: classical SVM with quantum kernel -- can find patterns invisible classically
\end{itemize}
\begin{verbatim}
from qiskit_machine_learning.kernels import FidelityQuantumKernel
from qiskit.circuit.library import ZZFeatureMap
feature_map = ZZFeatureMap(feature_dimension=2, reps=2)
kernel = FidelityQuantumKernel(feature_map=feature_map)
\end{verbatim}
\end{frame}

%%%%%%%%%%%%%%%%%%%%%%%%%%%%%%%%%%%%%%%%%%%%%%%%%%%%%%%%%%%
\begin{frame}[fragile]\frametitle{Mini-Project: Full Qiskit Pipeline}
\begin{itemize}
\item Task: Implement Grover's search for a 4-qubit problem (16 items, mark ``1010'')
\item Build the oracle using phase kickback
\item Apply Grover diffusion operator; run for optimal number of iterations
\item Test on Aer simulator; then submit to IBM Quantum hardware
\item Compare simulation vs hardware results
\item Report: circuit depth, qubit count, fidelity on hardware
\end{itemize}
\end{frame}

%%%%%%%%%%%%%%%%%%%%%%%%%%%%%%%%%%%%%%%%%%%%%%%%%%%%%%%%%%%
\begin{frame}[fragile]\frametitle{Qiskit Cheatsheet: Quick Reference}
\begin{itemize}
\item \textbf{Create}: \texttt{QuantumCircuit(n, m)}
\item \textbf{Gates}: \texttt{.h(q)}, \texttt{.x(q)}, \texttt{.cx(c,t)}, \texttt{.ry(theta,q)}
\item \textbf{Measure}: \texttt{.measure(qubits, cbits)}
\item \textbf{Draw}: \texttt{.draw('mpl')}
\item \textbf{Simulate}: \texttt{AerSimulator().run(transpile(qc, sim), shots=1024)}
\item \textbf{State}: \texttt{Statevector(qc)}
\item \textbf{Transpile}: \texttt{transpile(qc, backend, optimization\_level=3)}
\item \textbf{Hardware}: \texttt{QiskitRuntimeService(); SamplerV2()}
\end{itemize}
\end{frame}

%%%%%%%%%%%%%%%%%%%%%%%%%%%%%%%%%%%%%%%%%%%%%%%%%%%%%%%%%%%
\begin{frame}[fragile]\frametitle{Module 11 Summary}
\begin{itemize}
\item Qiskit is the primary Python framework for quantum programming
\item QuantumCircuit: build, visualize, and run circuits on simulators and hardware
\item Aer simulators: StatevectorSimulator (exact), AerSimulator (sampled, with noise)
\item Parameterized circuits: essential for VQE, QAOA, and QML
\item Runtime Primitives (Sampler/Estimator): unified hardware/simulator interface in Qiskit 1.x
\item Full pipeline: circuit $\to$ transpile $\to$ simulate/submit $\to$ analyze results
\item Practice: the best way to learn quantum programming is to run circuits!
\end{itemize}
\end{frame}
